% Vector redraw of bottom_slide.jpg, preserving its 960 x 545 pixel canvas.
% Regenerate from this directory:
% xelatex -no-pdf -interaction=nonstopmode -halt-on-error bottom-slide.tex
% dvisvgm --no-fonts --exact --bbox=papersize --output=bottom-slide.svg bottom-slide.xdv
\def\pgfsysdriver{pgfsys-dvisvgm.def}
\documentclass[tikz,border=0pt]{standalone}
\usepackage[UTF8,fontset=fandol]{ctex}
\begin{document}
\begin{tikzpicture}[x=0.75bp,y=-0.75bp,line cap=round,line join=round]
  \path[use as bounding box] (0,0) rectangle (960,545);
  \clip (0,0) rectangle (960,545);
  \fill[white] (0,0) rectangle (960,545);

  % A general partition into distinct parts; ellipses denote omitted dots/rows.
  \foreach \x/\y in {
    160/36,225/36,390/36,455/36,620/36,774/36,844/36,
    160/106,225/106,390/106,455/106,620/106,774/106,
    160/260,225/260,390/260,455/260,620/260,
    160/410,225/410,390/410
  } {
    \fill (\x,\y) circle[radius=6bp];
  }
  % Horizontal ellipses in each visible row.
  \foreach \x/\y in {
    305/36,535/36,700/36,
    305/106,535/106,700/106,
    305/260,535/260,
    305/410
  } {
    \foreach \offset in {-20,0,20} {
      \fill ({\x+\offset},\y) circle[radius=1.2bp];
    }
  }
  % Vertical ellipses between the second and third visible rows.
  \foreach \x in {160,225,390,455,620} {
    \foreach \y in {166,182,198} {
      \fill (\x,\y) circle[radius=1.2bp];
    }
  }
  % Vertical ellipses before the bottom row.
  \foreach \x in {160,225,390} {
    \foreach \y in {318,334,350} {
      \fill (\x,\y) circle[radius=1.2bp];
    }
  }

  % The slide joins the rightmost dots at 45 degrees.
  \draw[line width=1.25bp] (620,260) -- (844,36);
  \node[inner sep=0pt,font=\kaishu\fontsize{24bp}{28bp}\selectfont] at (766,177) {坡};

  % A brace identifies the bottom row.
  \draw[line width=1.5bp]
    (151,454) .. controls (155,463) and (162,467) .. (174,467)
    -- (248,467) .. controls (259,467) and (265,470) .. (270,477)
    .. controls (275,470) and (281,467) .. (292,467)
    -- (363,467) .. controls (378,467) and (384,462) .. (388,454);
  \node[inner sep=0pt,font=\kaishu\fontsize{24bp}{28bp}\selectfont] at (270,514) {底};
\end{tikzpicture}
\end{document}
